\documentclass[type=bachelor,fontset=fandol]{thuthesis}
% 选项：
%   type=[bachelor|master|doctor|postdoctor], % 必选
%   fontset=[fandol|custom],                  % 字体选择
%   secret,                                   % 可选
%   pifootnote,                               % 可选（建议打开）
%   openany|openright,                        % 可选，基本不用
%   arial,                                    % 可选，基本不用
%   arialtoc,                                 % 可选，基本不用
%   arialtitle                                % 可选，基本不用

% **生僻字乱码的解决方法**
% 先解释一下原因: 自带的 fandol 字体不包含生僻字，常用的中易字体因为版权
% 原因不能在 overleaf 发布。
% 解决方法: 找一台 windows 电脑，从 C:\Windows\Fonts 里把
%   simsun.ttc simhei.ttf simkai.ttf simhei.ttf
% 复制出来，上传到 fonts 文件夹，再设置 typeset=custom 即可。


% 所有其它可能用到的包都统一放到这里了，可以根据自己的实际添加或者删除。
\usepackage{thuthesis}

% 定义所有的图片文件在 figures 子目录下
\graphicspath{{figures/}}

% 可以在这里修改配置文件中的定义。导言区可以使用中文。
% \def\myname{薛瑞尼}

\begin{document}

%%% 封面部分
\frontmatter
\thusetup{
  ctitle={深度神经网络性能预测模型},
  cdegree={工学学士},
  cdepartment={计算机科学与技术系},
  cmajor={计算机科学与技术},
  cauthor={马子轩},
  csupervisor={陆游游},
}

\begin{cabstract}
    随着机器学习的快速发展，深度神经网络在计算机视觉、自然语言处理等领域有着广泛的应用。深度神经网络规模越来越大，训练时间越来越长，需要的计算资源也越来越多。因此，针对神经网络应用使用的计算资源的预算控制、训练过程的配置选择、训练集群的资源分配都成为了迫在眉睫的任务，而这些任务的完成，需要对深度神经网络的运行时间有准确的预测。

    传统的预测方法存在各类问题，集中体现在：对复杂模型的支持不够好，没有考虑操作间的依赖关系。无法支持非线性操作，对资源利用率不足的场景无法预测等方面。
    
    针对传统方法的两大问题，本文提出了基于“预测-调度”框架的深度神经网络性能预测模型。通过对深度神经网络结构进行静态分析，预测神经网络的运行时间。模型主要分为性能模型和调度模拟两部分。其中，性能模型部分分别对各个操作的运行时间进行预测。调度模拟模拟数据流图的运行过程，得到操作的执行顺序和并发情况。两者结合得到准确的模型运行时间。
    
    和传统预测方法相比，我们的实现能够更好地解决复杂的模型调度过程以及运行过程中资源利用不充分导致的性能预测不准确的问题。我们的模型平均准确度能达到80\%以上，和TensorFlow提供的基于传统预测方式的预测工具相比，我们的模型在最优测试中能提高100倍以上。

\end{cabstract}

\ckeywords{深度神经网络, 机器学习, 性能预测}

\begin{eabstract}
    With the rapid development of machine learning, deep neural networks(DNNs) have been widely used in computer vision, natural language processing and other fields. DNNs are getting larger, and training time is getting longer, more computing resources are needed. Therefore, the budget control of computational resources for DNN applications, the configuration selection of training process and the resource allocation of training cluster are all urgent tasks. The completion of these tasks requires accurate prediction model of the running time of DNN.

    There are various problems with traditional prediction models, which are: the support of complex models is not good enough, and the dependency between operations is not considered. Nonlinear operations cannot be supported, and insufficient resource utilization cannot be predicted.

    Aiming at the two problems of traditional methods, this paper proposes a DNN performance prediction model based on the framework of "prediction - scheduling". The running time of DNN is predicted by static analysis of its structure. The model is divided into two parts: performance model and scheduling simulator. In the performance model section, the running time of each operation is predicted respectively. Scheduling simulator simulates the running process of data flow graph, and obtains the execution order and concurrency of operations. The combination of the two results gives us accurate running time.

    Compared with traditional prediction model, our implementation can better solve the problem of inaccurate performance prediction caused by the complex model and the insufficient utilization of resources. Our model can achieve an average accuracy of more than 80\%, which is more than 100 times better than the prediction tools provided by TensorFlow based on traditional prediction model.
\end{eabstract}

\ekeywords{DNN, ML, Performance}

% 如果使用授权说明扫描页，将可选参数中指定为扫描得到的 PDF 文件名，例如：
% \makecover[scan-auth.pdf]
\makecover

%% 目录
\tableofcontents

%% 符号对照表
\begin{denotation}[3cm]
    \item[ML] 机器学习 (Machine learning)
    \item[DNN] 深度神经网络 (Deep neural network)
    \item[CNN] 卷积神经网络 (Convolutional neural network)
    \item[CPU] 中央处理器 (Central processing unit)
    \item[GPU] 图形处理器 (Graphics processing unit)
\end{denotation}



%%% 正文部分
\mainmatter
\chapter{介绍}
\label{cha:intro}

\section{动机}
\label{sec:motivation}
    随着机器学习的发展，深度神经网络有了越来越广泛的应用，如在机器视觉，自然语言处理等领域。于此同时，深度神经网络的深度和规模都越来越大，因此，深度神经网络的运行需要{\bfseries 更大规模的运算设备以及更长的运行时间}。那么，选择哪种设备或知道在固定设备上，模型的运行时间，就成了一个重要的问题。

    运行机器学习应用有两种不同的模式。在本地设备运行或在云平台运行。通常，企业会有自己的设备集群，并配备相应的管理方式。这种情况下，计算资源成为了企业内部争抢的资源，这时如何进行任务层面的调度，就成为了一个非常重要的问题。而更准确的离线性能预测模型，能够给企业提供更好的调度效果，提高工作效率。
    
    另一方面，当资金比较缺乏，不能够自己组建设备，或需要的计算资源规模太大(如Facebook 2018年的工作一样，使用256块GPU在1小时内训练ImageNet)。这些情况下，在本地运行都不现实，因此，使用云计算平台成为了非常合适的解决方案。
    
    云计算平台上通常提供了现成的机器学习应用框架，如AWS Machine Learning, Azure Machine Learning等，他们提供了从现有模型导入到云计算平台运行的功能。而另一种方式是直接租用实例（如AWS EC2等）搭建机器学习框架进行运行。无论使用哪种方式，现有的云计算平台都是按照时间收费的。因此，预测模型在不同设备上的运行时间，能够帮助用户更好的规划资金，选择配置。
    
    而另一方面，现有的机器学习框架，如tensorflow等，在集群模式下，通常使用数据并行方式进行加速。而在模型并行模式下，由于缺少操作性能数据模型，不能得到每个操作较为准确的运行时间，限制了模型并行的细粒度调度。因此，在进行细粒度性能预测的过程中，模型能够给机器学习框架提供细粒度调度建议，提高模型并行的运行速度。
    
    综上所述，深度神经网络性能预测模型无论是对工业界还是学术界都很重要。

\section{相关工作}
\label{sec:related}

\subsection{Tensorflow}
    深度学习应用通常运行在深度学习框架上，
\chapter{中华人民共和国}
\label{cha:china}

\section{其它例子}
\label{sec:other}

在第~\ref{cha:intro} 章中我们学习了贝叶斯公式~(\ref{equ:chap1:bayes})，这里我们复
习一下：
\begin{equation}
\label{equ:chap2:bayes}
p(y|\mathbf{x}) = \frac{p(\mathbf{x},y)}{p(\mathbf{x})}=
\frac{p(\mathbf{x}|y)p(y)}{p(\mathbf{x})}
\end{equation}

\subsection{绘图}
\label{sec:draw}

本模板不再预先装载任何绘图包（如 \pkg{pstricks，pgf} 等），完全由用户来决定。
个人觉得 \pkg{pgf} 不错，不依赖于 Postscript。此外还有很多针对 \LaTeX{} 的
 GUI 作图工具，如 XFig(jFig), WinFig, Tpx, Ipe, Dia, Inkscape, LaTeXPiX,
jPicEdt, jaxdraw 等等。

\subsection{插图}
\label{sec:graphs}

强烈推荐《\LaTeXe\ 插图指南》！关于子图形的使用细节请参看 \pkg{subcaption} 宏包的说明文档。

\subsubsection{一个图形}
\label{sec:onefig}
一般图形都是处在浮动环境中。之所以称为浮动是指最终排版效果图形的位置不一定与源文
件中的位置对应\footnote{This is not a bug, but a feature of \LaTeX!}，这也是刚使
用 \LaTeX{} 同学可能遇到的问题。如果要强制固定浮动图形的位置，请使用 \pkg{float} 宏包，
它提供了 \texttt{[H]} 参数，比如图~\ref{fig:xfig1}。
\begin{figure}[H] % use float package if you want it here
  \centering
  \includegraphics{thu-whole-logo}
  \caption{利用 Xfig 制图}
  \label{fig:xfig1}
\end{figure}

大学之道，在明明德，在亲民，在止于至善。知止而后有定；定而后能静；静而后能安；安
而后能虑；虑而后能得。物有本末，事有终始。知所先后，则近道矣。古之欲明明德于天
下者，先治其国；欲治其国者，先齐其家；欲齐其家者，先修其身；欲修其身者，先正其心；
欲正其心者，先诚其意；欲诚其意者，先致其知；致知在格物。物格而后知至；知至而后
意诚；意诚而后心正；心正而后身 修；身修而后家齐；家齐而后国治；国治而后天下
平。自天子以至于庶人，壹是皆以修身为本。其本乱而未治者 否矣。其所厚者薄，而其所
薄者厚，未之有也！

\hfill —— 《大学》


\subsubsection{多个图形}
\label{sec:multifig}

如果多个图形相互独立，并不共用一个图形计数器，那么
用 \texttt{minipage} 或者\texttt{parbox} 就可以。否则，请参看
图~\ref{fig:big1-subcaptionbox}，它包含两个小图，分别是图~\ref{fig:subfig1}和
图~\ref{fig:subfig2}。推荐使用\cs{subcaptionbox}，因为可以像
图~\ref{fig:big1-subcaptionbox} 那样对齐子图的标题，也可以使
用\pkg{subcaption}宏包的\cs{subcaption}（放在 minipage中，用法同\cs{caption}）
或是 \pkg{subfigure} 、 \pkg{subtable}环境，像图~\ref{fig:big1-subfigure}，
不要再用 \cs{subfloat}、\cs{subfigure} 和 \cs{subtable}。

\begin{figure}[h]
  \centering%
  \subcaptionbox{第一个小图形\label{fig:subfig1}}[3cm] %标题的长度，超过则会换行，如下一个小图。
    {\includegraphics[height=3cm]{thu-fig-logo}}%
  \hspace{4em}%
  \subcaptionbox{第二个小图形，注意这个图略矮些。如果标题很长的话，它会自动换行\label{fig:subfig2}}
      {\includegraphics[height=2cm]{thu-text-logo}}
  \caption{包含子图形的大图形(subcaptionbox示例)}
  \label{fig:big1-subcaptionbox}
\end{figure}
\begin{figure}[h]
  \centering%
  \begin{subfigure}{3cm}
    \includegraphics[height=3cm]{thu-fig-logo}
    \caption{第一个小图形}
  \end{subfigure}%
  \hspace{4em}%
  \begin{subfigure}{0.5\textwidth}
    \includegraphics[height=2cm]{thu-text-logo}
    \caption{第二个小图形，注意这个图略矮些。subfigure中同一行的子图在顶端对齐。}
  \end{subfigure}
  \caption{包含子图形的大图形(subfigure示例)}
  \label{fig:big1-subfigure}
\end{figure}

古之学者必有师。师者，所以传道受业解惑也。人非生而知之者，孰能无惑？惑而不从师，
其为惑也，终不解矣。生乎吾前，其闻道也固先乎吾，吾从而师之；生乎吾後，其闻道也亦
先乎吾，吾从而师之。吾师道也，夫庸知其年之先後生於吾乎！是故无贵无贱无长无少，道
之所存，师之所存也。

嗟乎！师道之不传也久矣，欲人之无惑也难矣。古之圣人，其出人也远矣，犹且从师而问焉；
今之众人，其下圣人也亦远矣，而耻学於师。是故圣益圣，愚益愚。圣人之所以为圣，愚
人之所以为愚，其皆出於此乎？爱其子，择师而教之，於其身也，则耻师焉，惑焉。彼童子
之师，授之书而习其句读者，非吾所谓传其道、解其惑者也。句读之不知，惑之不解，或师
焉，或不焉，小学而大遗，吾未见其明也。巫医、乐师、百工之人不耻相师，  士大夫之族
曰“师”曰“弟子”之云者，则群聚而笑之。问之，则曰：彼与彼年相若也，道相似也，位
卑则足羞，官盛则近谀。呜呼！师道之不复，可知矣。巫医、乐师、百工之人。吾子不齿，
今其智乃反不能及，其可怪也欤！圣人无常师。孔子师郯子、苌子、师襄、老聃。郯子之徒，
其贤不及孔子。孔子曰：“三人行，必有我师。”是故弟子不必不如师，师不必贤於弟子。
闻道有先後，术业有专攻，如是而已。

如果要把编号的两个图形并排，那么小页就非常有用了：
\begin{figure}
\begin{minipage}{0.48\textwidth}
  \centering
  \includegraphics[height=2cm]{thu-whole-logo}
  \caption{并排第一个图}
  \label{fig:parallel1}
\end{minipage}\hfill
\begin{minipage}{0.48\textwidth}
  \centering
  \includegraphics[height=2cm]{thu-whole-logo}
  \caption{并排第二个图}
  \label{fig:parallel2}
\end{minipage}
\end{figure}

李氏子蟠，年十七，好古文、六艺，经传皆通习之，不拘於时，学於余。余嘉其能行古
道，作师说以贻之。

\hfill —— 韩愈（唐）



%%% 其它部分
\backmatter

%% 本科生要这几个索引，研究生不要。选择性留下。
% 插图索引
\listoffigures
% 表格索引
\listoftables
% 公式索引
\listofequations


%% 参考文献
% 注意：至少需要引用一篇参考文献，否则下面两行可能引起编译错误。
% 如果不需要参考文献，请将下面两行删除或注释掉。
% \bibliographystyle{thuthesis-numerical}
% \bibliography{ref/refs}


%% 致谢
\begin{acknowledgement}
    感谢陆游游老师在毕设过程中对我的悉心指导。
    
    感谢冯冠宇、谢磊、郑帧三位学长在毕设期间对我的贴心帮助以及提出富有建设性的意见。
    
    感谢实验室全体师生对我的帮助。

\end{acknowledgement}


%% 附录
\begin{appendix}
\chapter{外文资料的调研阅读报告}
\title{TensorFlow}

\section{引言}
人工智能在当今社会中应用越来越广泛，机器学习的地位与日俱增。深度神经网络在计算机视觉、自然语言处理等领域都有广泛的应用。随着数据规模和网络规模的日益扩大，以及异构平台的广泛使用，机器学习框架的重要性与日俱增。现有的机器学习框架有很多，如TensorFlow、Caffe、Torch、Theano等。本文将对TensorFlow进行简单的介绍。

\section{历史}
2010年Google就开发了他们的第一代机器学习系统DistBelief。这一系统在Google的搜索、广告、翻译等多个服务中都有广泛的应用。后来在Jeff Dean等人在DistBelief的基础上进行优化，创造了TensorFlow。现在TensorFlow已经代替DistBelief成为Google主要使用的系统。

\section{介绍}
TensorFlow的名字来源于TensorFlow的运行过程。Tensor是张量的意思，可以理解为一个多维数组。TensorFlow的计算被组织成数据流图的形式，Tensor穿梭于图上的节点之间，就像流动一样，所以起名TensorFlow。

\section{为什么要使用TensorFlow}
在我们执行机器学习任务的时候，用户往往更关注的是模型的效果、如收敛时间、精确度、等信息，而不是模型该如何部署到大规模的集群，如何处理运行间的错误等。如果使用C++进行过神经网络的编写，你就会知道，一个仅仅几层的神经网络，从正向预测到反向传播，整个代码会有几百行之长。这是机器学习算法工程师不能接受的，严重阻碍了机器学习的发展。另一方面，机器学习任务的模型本身通常是不受平台限制的。比如同样的一个网络，可以运行在单机CPU，也可以执行在大规模异构集群上，而部署阶段甚至可以运行在边缘设备，如手机等平台上。这对数据和模型的变化可能仅仅在于一两个参数的变化，但是如果没有机器学习框架的帮助，需要重复在不同平台上实现多个不同版本。

Tensorflow给用户主要提供了两个功能。一个是提供给用户简单易用的编程接口，方便用户在使用的时候定义机器学习任务。另一个是帮助用户执行模型。在TensorFlow中，系统预先定义了大量操作，用户可以使用TensorFlow提供的基本原语，非常容易地创建机器学习任务，不需要关心过多的细节。同时TensorFlow提供了执行框架，相同的机器学习任务可以轻松的移植到不同的平台，从手机到大规模异构集群，仅仅需要少量的改动即可使用

\section{数据流图}
数据流图是一个古老的工具，在80年代就有人使用数据流图来进行并行计算。本世纪初，数据流图经历了历史的寒冬，正如人工智能所经历的一样。近些年，计算力的大幅提升、以及异构平台的大量使用，给了数据流图新的历史发展机遇。

数据流图的好处在于，能够将复杂的计算任务高度抽象，节点间的运行互不影响。在这个前提下，任务能够被比较好地调度到不同的设备上进行执行。

现在，大量最先进的数据处理系统，都使用数据流图作为编程模型，或数据的抽象模型。如Spark、Flink等。

\section{TensorFlow的编程模型}
TensorFlow中机器学习任务被组织成数据流图的形式。数据流图是一个有向无环图，其中图上的节点表示TensorFlow中定义好的一个运算操作，图上的边表示数据在节点之间的流动过程。特别地，在TensorFlow中，数据被以张量（Tensor）的形式存储。图上的边不仅表示数据的流动过程，还同时表示了操作之间的依赖关系。另外，TensorFlow的数据流图中还有另一种边，它没有数据经过，用来限制计算之间的顺序关系，这种边被称作控制依赖。

\section{操作与核函数}
在我们定义机器学习任务的时候，我们使用了TensorFlow中提供的各类操作。这些操作只是标记了一种运算方式。我们在执行计算任务的时候，操作可能会被映射到不同的实现上。这些实现被称作核函数。

一般情况下，核函数和操作以及设备进行绑定。一个操作可能对应多个核函数，每个核函数对应一种类型的设备。大部分的核函数使用Eigen实现，在GPU上很多核函数是用CuBlas和CuDNN函数实现的，因此TensorFlow能够在不同平台上都有比较不错的运行性能。

\section{会话}
会话是TensorFlow中的重要交互工具。会话主要提供两个基本功能：数据流图的创建和数据流图的执行接口。一般情况下，TensorFlow对每一个数据流图都会对应一个会话。在多次计算数据流图的时候，反复通过会话的启动命令进行运行。在会话的执行操作中，TensorFlow对数据流图进行处理、分配、以及运行过程。

\section{TensorFLow的实现}
TensorFlow的计算和客户端是分开的，用户定义计算任务是在客户端进行的。用户在客户端可以通过会话连接master，向master提交任务，传输数据流图。Master对数据流图进行处理，再发送给worker。Worker根据得到的计算任务分配到本地设备，分别执行对应的核函数。

其中，master在拿到任务后，先进行一次剪枝操作。在这步操作中，TensorFlow执行DFS算法，将不会被执行到的操作去除。之后，根据平台的配置情况将处理过的图切分，分配到不同的worker上。例如单机的时候，worker仅有一个，就会直接将图传送给worker进程。而在分布式部署的时候，就会将切分好的图分配到不同的worker上。Worker在拿到被分配的子图后，还会进行第二次切分，将节点切分到不同的设备上，再调用数据流图执行算法。

单机单设备执行过程，单机单GPU执行是最简单的情况。这时候客户端、master和worker都是本地的一个进程。这时模型最为简单，有一个队列用来存放待执行任务。每一个待执行任务都有一个依赖计数器。每当一个任务被执行完成的时候，就给他的出边对应的待执行操作的计数器减1。当依赖计数器为0的时候，这一个操作就可以被调度执行，调用当前设备上该操作对应的核函数。

单机多设备执行过程，在单机多设备的执行过程中，多了两个问题，一个是操作在哪个设备上执行。另一个是设备之间的数据依赖关系。通常，无论采用数据并行，还是模型并行。用户都需要提前指定操作被执行的设备。数据依赖上，多设备执行的时候，设备需要的数据可能存储在其他设备上，这时候需要插入send和receive节点，来进行数据的传输。Send/receive节点的插入有多种策略，并不一定在所有的数据依赖上都插入。其中允许的一种调度方式是定时以同步的方式在设备间传输数据。这种情况系统的可扩展性会更好。

分布式系统执行。分布式系统和多设备类似，只不过跨设备的数据依赖插入的send/receive节点以TCP或RDMA方式进行。

\section{总结}

TensorFlow是目前使用非常广泛的机器学习框架，我的毕设将在TensorFlow的基础上进行。目前，TensorFlow的调度过程可以比较容易地被模拟，我们在一系列实测操作数据的基础上就可以对TensorFlow的运行时间进行预测。我的毕设将在这一思想的基础上实现深度神经网络的性能预测。

\title{参考文献}

\begin{translationbib}
    \item ABADI M, BARHAM P, CHEN J, et al. Tensorflow: A system for large-scale machinelearning[C]//12thfUSENIXgSymposium on Operating Systems Design and Implementation(fOSDIg16). 2016: 265-283.
    \item ABADI M, AGARWAL A, BARHAM P, et al. TensorFlow: Large-scale machine learningon heterogeneous systems[EB/OL]. 2015.http://tensorflow.org/.
    \item BAHRAMPOUR S, RAMAKRISHNAN N, SCHOTT L, et al. Comparative study of deeplearning software frameworks[J]. arXiv preprint arXiv:1511.06435, 2015.
\end{translationbib}
\end{appendix}

%% 个人简历
\begin{resume}

  \resumeitem{个人简历}

  xxxx 年 xx 月 xx 日出生于 xx 省 xx 县。

  xxxx 年 9 月考入 xx 大学 xx 系 xx 专业，xxxx 年 7 月本科毕业并获得 xx 学士学位。

  xxxx 年 9 月免试进入 xx 大学 xx 系攻读 xx 学位至今。

  \researchitem{发表的学术论文} % 发表的和录用的合在一起

  % 1. 已经刊载的学术论文（本人是第一作者，或者导师为第一作者本人是第二作者）
  \begin{publications}
    \item Yang Y, Ren T L, Zhang L T, et al. Miniature microphone with silicon-
      based ferroelectric thin films. Integrated Ferroelectrics, 2003,
      52:229-235. (SCI 收录, 检索号:758FZ.)
    \item 杨轶, 张宁欣, 任天令, 等. 硅基铁电微声学器件中薄膜残余应力的研究. 中国机
      械工程, 2005, 16(14):1289-1291. (EI 收录, 检索号:0534931 2907.)
    \item 杨轶, 张宁欣, 任天令, 等. 集成铁电器件中的关键工艺研究. 仪器仪表学报,
      2003, 24(S4):192-193. (EI 源刊.)
  \end{publications}

  % 2. 尚未刊载，但已经接到正式录用函的学术论文（本人为第一作者，或者
  %    导师为第一作者本人是第二作者）。
  \begin{publications}[before=\publicationskip,after=\publicationskip]
    \item Yang Y, Ren T L, Zhu Y P, et al. PMUTs for handwriting recognition. In
      press. (已被 Integrated Ferroelectrics 录用. SCI 源刊.)
  \end{publications}

  % 3. 其他学术论文。可列出除上述两种情况以外的其他学术论文，但必须是
  %    已经刊载或者收到正式录用函的论文。
  \begin{publications}
    \item Wu X M, Yang Y, Cai J, et al. Measurements of ferroelectric MEMS
      microphones. Integrated Ferroelectrics, 2005, 69:417-429. (SCI 收录, 检索号
      :896KM)
    \item 贾泽, 杨轶, 陈兢, 等. 用于压电和电容微麦克风的体硅腐蚀相关研究. 压电与声
      光, 2006, 28(1):117-119. (EI 收录, 检索号:06129773469)
    \item 伍晓明, 杨轶, 张宁欣, 等. 基于MEMS技术的集成铁电硅微麦克风. 中国集成电路,
      2003, 53:59-61.
  \end{publications}

  \researchitem{研究成果} % 有就写，没有就删除
  \begin{achievements}
    \item 任天令, 杨轶, 朱一平, 等. 硅基铁电微声学传感器畴极化区域控制和电极连接的
      方法: 中国, CN1602118A. (中国专利公开号)
    \item Ren T L, Yang Y, Zhu Y P, et al. Piezoelectric micro acoustic sensor
      based on ferroelectric materials: USA, No.11/215, 102. (美国发明专利申请号)
  \end{achievements}

\end{resume}

\end{document}
